% Generated draft: requires mathematical and source review.
\documentclass[11pt]{article}
\usepackage[margin=25mm]{geometry}
\usepackage{fontspec}
\setmainfont{Latin Modern Roman}
\usepackage{amsmath,amssymb,mathtools,unicode-math}
\setmathfont{Latin Modern Math}
\usepackage{enumitem,needspace,xurl,hyperref}
\setlength{\parindent}{0pt}
\setlength{\parskip}{5pt}
\setlength{\emergencystretch}{3em}
\newcommand{\N}{\mathbb{N}}
\newcommand{\Z}{\mathbb{Z}}
\newcommand{\Q}{\mathbb{Q}}
\newcommand{\R}{\mathbb{R}}
\newcommand{\C}{\mathbb{C}}
\newcommand{\F}{\mathbb{F}}
\newcommand{\A}{\mathbb{A}}
\newcommand{\PP}{\mathbb{P}}
\newcommand{\E}{\mathbb{E}}
\newcommand{\eps}{\varepsilon}
\newcommand{\dd}{\,\mathrm{d}}
\begin{document}
\section*{UnsolvedMath ID: 30006059}
\textbf{Source problem code:} catalogue entry
\subsection{Definitions and mathematical statement}
For the formulation below, use the conventional smooth oriented knot setting: a knot is a smooth oriented embedded circle in the standard sphere \(S^3\). The supplied source does not explicitly specify ambient-space or orientation conventions; these conventions require confirmation for the intended statement.

For \(n\geq 1\), let \(B_n\) denote the Artin braid group with standard generators \(\sigma_1,\ldots,\sigma_{n-1}\). A positive braid is a braid admitting a word using only these generators, with no inverse generators. A knot is \emph{braid-positive} if it is isotopic to the closure of such a braid for some \(n\); the closure must have exactly one component. Allowing the empty word includes the identity braid in \(B_1\), whose closure is the unknot. The source does not explicitly address the empty-word convention.

Two oriented knots \(K_0,K_1\subset S^3\) are \emph{smoothly concordant} if there is a smooth properly embedded oriented annulus \(A\subset S^3\times[0,1]\) with oriented boundary
\[
\partial A=(-K_0)\times\{0\}\;\sqcup\;K_1\times\{1\},
\]
where \(-K_0\) denotes reversal of the knot orientation. They are \emph{isotopic} if there is a smooth ambient isotopy \(H_t:S^3\to S^3\), \(0\leq t\leq1\), with \(H_0=\operatorname{id}_{S^3}\) and \(H_1(K_0)=K_1\), respecting knot orientations under the stated convention.

\textbf{Question.} For every pair of braid-positive knots \(K_0,K_1\subset S^3\), does smooth concordance imply isotopy? Equivalently, under the conventions above, is
\[
\forall K_0,K_1\text{ braid-positive},\qquad
K_0\sim_{\mathrm{smooth\ concordance}}K_1
\ \Longrightarrow\ 
K_0\sim_{\mathrm{ambient\ isotopy}}K_1
\]
true?

\textbf{Supplied context and qualifications.} The excerpt states that braid-positive knots are fibered and strongly quasipositive, and presents this question as a special case of Baker's conjecture. Those properties are supplied background, not additional restrictions on the quantified class. The excerpt singles out closures of positive braids in \(B_3\) as a potentially accessible special case, citing classification of conjugacy classes of three-strand braids and earlier work. The target question nevertheless concerns positive braids on any number of strands. The excerpt also introduces \(\upsilon(K)=\Upsilon_K(1)\), described there as a smooth concordance invariant from knot Floer homology defined by Ozsv\'ath--Stipsicz--Szab\'o, and \(g(K)\), the Seifert genus. Neither invariant appears in the target assertion. The excerpt announces a corollary about positive three-braid knots but does not supply its statement; no such corollary is incorporated here.
\subsection{Short English statement}
The question asks whether two knots that each arise as the closure of a positive braid must represent the same knot type whenever they cobound a smooth concordance annulus. It concerns all braid-positive knots, with positive three-braid knots highlighted only as a special case. The ambient-space, orientation, and empty-word conventions above are explicit provisional conventions requiring review. The supplied registry labels the question open; that status is not independently verified here.
\subsection{Sources}
\begin{itemize}
\item Local UnsolvedMath registry, supplied complete source record: permanent ID \(30006059\), title \emph{Concordance Rigidity of Braid-Positive Knots}, exact source problem code \texttt{catalogue entry}, in the Oberwolfach Reports Open Problems collection. The record gives proposed year \(2024\) and no named proposer.
\item Supplied report citation: \emph{MATRIX-MFO Tandem Workshop: Invariants in Low-Dimensional Topology: Combinatorics, Geometry, and Computation} (2025), Oberwolfach Reports, pp.\ 2511--2550, DOI \texttt{10.4171/owr/2024/43}. \url{https://doi.org/10.4171/owr/2024/43}. The supplied literature-source label is \emph{Oberwolfach Report 43/2024}. The citation's 2025 date and the report's 2024 designation are retained as supplied.
\item Supplied literature source: Krishna--Morton, \emph{Twist positivity, L-space knots, and concordance}. \url{https://doi.org/10.1007/s00029-024-00996-6}.
\item Literature assessment dated exactly \texttt{2026-08-22}: the supplied curation labels the problem open and states that the universal concordance--isotopy question remains posed, while positive three-braid results provide invariants and infinite distinct concordance families without settling the universal assertion. This is supplied dated curation, not an independently verified literature assessment. The background expressly requests independent verification of time-sensitive status and cited claims.
\item The supplied registry marks the statement \texttt{corrected\_verified} and reports that the official Oberwolfach report was checked, the question isolated from surrounding exposition, and notation repaired. These are source-verification claims supplied by the registry, not verification performed for this draft.
\item The excerpt refers to internal references [7], [8], [9], [10], and [13], including [13, Proposition 3.2], but supplies no corresponding bibliography entries. It attributes the invariant \(\Upsilon\) to Ozsv\'ath--Stipsicz--Szab\'o and mentions Baker's conjecture without a complete citation or conjecture statement.
\end{itemize}
\end{document}
